

\section{Uniform C1: the bounded-reciprocal repair and its scope}

The preserved uniform-source theorem explicitly assumes $e$ infinitesimal.
The preserved contact--conductor theorem explicitly assumes $e$ unbounded.
This note proves the missing unit-$e$ case by extending the latter proof
to bounded $\Delta=e^{-1}$ and bounded testing $w$. The extension uses
the same exact complete parent, the same retained packets, and the same
unlocalized conductor. The new work is a weight estimate that remains
strict when both $\delta$ and $w$ have weight zero, followed by an
exhaustive testing-chart join. An old all-$e$ status assertion is not a
premise.

\begin{theorem}[Uniform C1, bounded reciprocal $e$]
\label{thm:bounded-delta}
Let the coefficient and testing fields be independent real closed
extensions of $\mathbb R$. Use the original canonical uniform C1 source
defined below, with $k,f,g>0$ infinitesimal, $e=kf/g>0$, and
$\Delta=1/e$ bounded. The shapes $a,t>0$ are bounded and satisfy the
original strict gaps $M=R-t>0$, $J=T-a>0$. Retain every original
coefficient sign, including $D_{\rm e},\beta\geq0$, and every literal
coefficient zero. Then the complete parent is tensor-SOS at every
independent projective measurement pair and every real or complex
testing state.
\end{theorem}

Together with the preserved infinitesimal-$e$ theorem, this proves the
whole positive-$e$ uniform source. It does not prove an unlisted source
chart or close any remaining C14 or weak-C2 row. We use the already
admitted separated bounded strict-Gram lemma: a finite bounded separated
Gram with strictly positive real residue gives tensor-SOS. No inversion
of a mixed element and no comparison between an independent coefficient
and testing rate enters the proof.

\section{Original source and complete parent}

Put $\kappa=k^2$, $\eta=f^2$, $\zeta=g^2$, $K=1+k^2$, and
\[
 V_c=\kappa\eta\zeta,\quad L_c=1-a\eta,\quad
 H_c=1+(1-a)\eta,\quad Q_c=1-\kappa\zeta+aV_c,
\]
\[
 R=1+k^2(1-a\eta),\qquad T=1+k^2-k^2\zeta t,
 \qquad M=R-t>0,\quad J=T-a>0.
\]
The native coquotients are defined by the following short polynomials:
\begin{align*}
 P_4={}&1+\kappa+\eta\kappa(a-\kappa-1)-\zeta\kappa t
 +V_c[-\zeta\kappa t^2+2\kappa t+a(1-t)+t],\\
 P_\alpha={}&1+\kappa-a\eta\kappa
 +\zeta[(1+\kappa)-(2+\kappa)t]
 +V_c[-\zeta t^2+a(t-1)+t],\\
 P_d={}&(1+\kappa)^2-a\eta\kappa^2
 +\zeta\kappa[(\kappa+1)-(\kappa+2)t]\\
 &+V_c[-\zeta\kappa t^2+\kappa a(t-1)+\kappa t+\kappa+1],\\
 c_2={}&2K(1-t\zeta)L_c(1+V_c)H_cQ_c,\\
 c_b={}&2KL_c[1+(1-t)\zeta]H_cQ_c,\\
 c_3={}&2KL_cH_cQ_c(1-\kappa\eta+tV_c),\\
 c_5={}&2KL_cH_cQ_c(1+V_c),\\
 c_4={}&2Q_cP_4,\qquad c_\alpha=2L_cH_cP_\alpha,
 \qquad c_d=4L_cQ_cP_d.
\end{align*}
Their residues are $2,2,2,2,2,2,4$. The original source is
\begin{gather}
 \lambda=n=0,\quad C=2,\quad
 d_{\rm B}=\frac{2a\eta c_d}{c_3},\quad
 m=\frac{2ae^2c_2}{Mc_3},\quad
 A=\frac{2a\eta^2c_4}{c_3},\quad
 \alpha=\frac{2a\Delta^2k^6\eta^2c_\alpha}{c_3},\notag\\
 D_{\rm e}+\frac{d_{\rm B}}2=\frac{2ac_5}{Jc_3},\qquad
 \beta+\frac{d_{\rm B}}2=\frac{2ae^2c_b}{tc_3},\qquad
 \Lambda=\frac{RTc_3}{2a}.                                      \tag{1}
\end{gather}
The original event $D_{\rm e}$ is distinct from the testing factor
$D=1+u^2$. For a testing state $(W,X,Y,Z)$ put
\[
 D=1+u^2,\quad D_v=1+v^2,\quad B=X+uZ,\quad Q_A=W+uY,
 \quad r=uW-Y,
\]
\[
 r_2=vQ_A-B,\qquad r_5=r-v(Z-uX).
\]
The original full parent on this source is
\begin{align*}
 F={}&4mr_2^2+8D_vr^2+4ADD_vY^2+(4D_{\rm e}+2d_{\rm B})r_5^2\\
 &+(4\alpha+2d_{\rm B})D(vY-Z)^2
 +(4\beta+2d_{\rm B})D_vB^2-2d_{\rm B}DD_vZ^2.
\end{align*}
With $H_2=4Tc_2$, $H_b=4Tc_b$, $H_3=4Rc_3$, $H_5=4Rc_5$,
equation (1) gives exactly
\begin{align}
 \Phi:=\Delta^2\Lambda F={}&H_2\frac RM r_2^2
 +H_b\frac RtD_vB^2+\Delta^2H_3\frac TaD_vr^2
 +\Delta^2H_5\frac TJr_5^2\notag\\
 &+4RT\Delta^2\eta^2c_4DD_vY^2
 +4RT\Delta^4k^6\eta^2c_\alpha D(vY-Z)^2\notag\\
 &+2RT\Delta^2\eta c_dD\{(vY-Z)^2-D_vZ^2\}.             \tag{2}
\end{align}

\section{The \texorpdfstring{$Z$}{Z}-unit chart with bounded \texorpdfstring{$W$}{W}}

Assume $u,v,X,Y$ are infinitesimal, $Z$ is a unit, and the state has
bounded coordinates. Normalize $Z=1$ in the testing field. The new
argument allows arbitrary bounded $W$. Put
\[
 p=t/R,\quad h_s=a/T,\quad q_A=v(W+uY),\quad q_B=v(1-uX),
 \quad N=B-pq_A,\quad M_B=r-h_sq_B.
\]
Subtract only the four nonnegative packets
\[
 \Pi_1=H_2\frac tM N^2,\quad \Pi_2=H_b\frac MtD_vN^2,
 \quad \Pi_3=\Delta^2H_3\frac JaD_vM_B^2,
 \quad \Pi_4=\Delta^2H_5\frac aJM_B^2.
\]
No boundedness of their gap ratios is assumed. The exact remainder is
\begin{align}
 G:=\Phi-\sum_i\Pi_i={}&H_2\{N^2-2q_AN+(1-p)q_A^2\}
 +H_bD_v\{N^2+2q_AN+pq_A^2\}\notag\\
 &+\Delta^2H_3D_v\{M_B^2+2q_BM_B+h_sq_B^2\}\notag\\
 &+\Delta^2H_5\{M_B^2-2q_BM_B+(1-h_s)q_B^2\}
 +\mathcal W,                                                   \tag{3}
\end{align}
where $\mathcal W$ is the last three terms of (2), with $Z=1$.
Use the pure-testing coordinate $w=(W+uY)/D$. Its full inverse is
\begin{gather}
 X=N+pDvw-u,\qquad q_B=v(D-uN-puDvw),\notag\\
 W=w+\frac{u(M_B+h_sq_B)}D,\qquad
 Y=uw-\frac{M_B+h_sq_B}D.                               \tag{4}
\end{gather}
The only testing denominator is the positive unit $D$.

If the pure-testing signed core is nonnegative, its product with
$d_{\rm B}>0$ is already a separated square and all event packets are
nonnegative. Otherwise the sign symmetries allow $u,v>0$. Explicitly,
changing $u$ to $-u$ and $(Y,Z)$ to $(-Y,-Z)$, or changing $v$ to
$-v$ and $(X,Z)$ to $(-X,-Z)$, preserves every event square and the
core. The negative-core branch then forces $w>0$: completing the two
positive core squares leaves
\[
 \frac{uvDD_v}{\mathcal D}
 \{D u v\,w^2-2DD_vw+uv\},\qquad
 \mathcal D=u^2v^4+u^2v^2+u^2+v^4+2v^2+1>0,
\]
which is nonnegative for $w\leq0$. The eliminated matrix is
\[
 \begin{pmatrix}1+v^2+u^2v^2&uv\\uv&(D+v^2)/D\end{pmatrix}\succ0.
\]
Thus all radicals used below have nonnegative pure-field factors.
Literal angle axes already have a nonnegative core.

\subsection{Polynomial source change and exact scalar}

Define pure-source quantities
\[
 \delta=\Delta R/K,\quad b^2=\eta K/R,\quad
 Q=K+ak^2b^2,\quad T_n=K^2-t\delta^2k^4Qb^2.
\]
Here $\delta$ is positive bounded, $k,b$ are infinitesimal, and
\begin{equation}
 \eta=Kb^2/Q,\quad \Delta=\delta Q/K,\quad
 \zeta=\delta^2k^2Qb^2/K,\quad R=K^2/Q,\quad T=T_n/K.       \tag{5}
\end{equation}
Every denominator in (5) is a positive source unit. After this
substitution put $C_i=K^2Q^3c_i$. These are polynomials: each native
monomial $c k^jf^{2r}g^{2s}a^\alpha t^\tau$ maps to
\[
 c\delta^{2s}k^{j+2s}b^{2(r+s)}a^\alpha t^\tau
 K^{2+r-s}Q^{3-r+s},\qquad -2\leq r-s\leq3.
\]
No negative exponent or shape inverse occurs. Put
\[
 V_n=K^2-tQuvw,\quad Y_n=KT_nuw-avV_n,\quad B_0=vY_n-KT_n,
\]
\[
 A_c=C_5T_n+aK(C_3D_v-C_5),\qquad \mathcal C=K^6Q^4T_n.
\]
Use different letters for the conductor monomials and the event weights:
\[
 A_m=wv,\quad B_m=\delta v,\quad C_m=\delta b^2uw,
 \quad D_0=\delta^2k^3b^2.
\]
Direct substitution of (4) into (3), retaining every term, gives
\begin{align}
 F_0:=(\mathcal CG)|_{N=M_B=0}={}&
 4KQT_n^2[C_2(K^2-tQ)+C_bD_vtQ]D^2 A_m^2\notag\\
 &+4Q^2A_cD^2V_n^2B_m^2\notag\\
 &+4K^3C_4DD_v(KT_nC_m-ab^2B_mV_n)^2\notag\\
 &+4KQ^2C_\alpha DB_0^2D_0^2\notag\\
 &+2K^2QC_dD\{b^2B_m^2(Y_n^2-K^2T_n^2+2aKT_nV_n)
                         -2K^2T_n^2B_mC_m\}.             \tag{6}
\end{align}
Define the leading polynomial and its complete error by
\[
 P=D(A_m^2+B_m^2)+K(C_m^2+D_0^2)-2K^2B_mC_m,
 \qquad E_0=F_0-8K^2D\mathcal CP.
\]
The full first coefficients on $(N,\delta M_B)$ are
\begin{align}
 S_B={}&-8aK^2b^4C_4D_vV_n
       +8\delta^2Q^2k^6b^4C_\alpha B_0+4KQb^2C_dB_0,\notag\\
 S_C={}&8K^3T_nb^2C_4D_v,\notag\\
 L_1={}&A_m[8K^3QT_n^2(C_bD_v-C_2)D]\notag\\
 &+B_m[-8\delta K^2Q^2uD\,vV_nA_c+\delta K^3auvS_B]
   +C_m[\delta K^3auvS_C],\notag\\
 L_2={}&B_m[8K^2Q^2T_nDV_n(C_3D_v-C_5)-K^2T_nS_B]
          -C_mK^2T_nS_C.                               \tag{7}
\end{align}
In particular $E_0\in(A_m,B_m,C_m)^2+(D_0^2)$ and
$L_1,L_2\in(A_m,B_m,C_m)$ without ideal division.

\subsection{Exact contact and the unlocalized conductor}

Work in $\mathcal R=\Rc[u,w,v]$ and put
\[
 \begin{gathered}
 \mathfrak J=(u-k,w-\delta k,v-\delta k^2b^2),\\
 x=B_m-C_m,\quad y=A_m-kC_m,\quad z=uB_m-kC_m,
 \quad h=kC_m-D_0.
 \end{gathered}
\]
The original source contact is
\[
 u=k,\quad v=\Delta k^2\eta,\quad
 (W,X,Y,Z)=(\Delta k,-k(1-t\zeta),\Delta k^2L_c,1).
\]
It becomes $(u,w,v,N,M_B)=(k,\delta k,\delta k^2b^2,0,0)$.
For completeness, its exact contact and derivative identities follow
from the six native balances
\[
 c_b-c_2=\zeta c_3,\qquad c_3-c_5=-\kappa\eta c_b,
 \qquad c_2=(1-t\zeta)c_5,
\]
\begin{align*}
 0={}&2\kappa c_2+2(1+V_c)L_cc_4
         -2\kappa^2\zeta Q_cc_\alpha-Q_cc_d,\\
 0={}&\kappa L_cc_2+c_5+(1+V_c)L_c^2c_4+\kappa Q_c^2c_\alpha
         +\tfrac12(\kappa\zeta L_c^2-2L_c-\eta)c_d,\\
 0={}&\kappa Rc_2+t\kappa V_cc_b+aV_cc_3+Tc_5
       +KV_cL_c^2c_4-\kappa^2\zeta KQ_cL_cc_\alpha
       -\tfrac K2(Q_cL_c+\eta)c_d.
\end{align*}
They follow by multiplying the short $P_4,P_\alpha,P_d$ polynomials.
Substitution into the derivatives of (2) proves that its value and all
first derivatives vanish at the displayed contact. The four packets
also vanish there to second order. Consequently
\[
 E_0\in\mathfrak J^2,\qquad L_1,L_2\in\mathfrak J.       \tag{8}
\]
Here $P\in\mathfrak J^2$ follows also from its exact square identity
below. No asymptotic contact replacement is made.

Adjoin only the separated radicals
\[
 a_1=\sqrt{wv},\quad a_2=\sqrt\delta\sqrt{uv},\quad
 a_3=b\sqrt{\delta k}\sqrt{uw},\quad
 a_4=\delta k^{3/2}b,\qquad q=a_1a_2-a_3a_4.
\]
Only $q^2\in\mathcal R$ is needed in the polynomial lift. Direct
multiplication gives
\begin{align}
 C_m(u-k)&=z-ux,\notag\\
 C_m(w-\delta k)&=\delta y-wx,\notag\\
 C_m(v-\delta k^2b^2)&=b^2(wz+\delta ky-kwx),\notag\\
 r_*:=y+z-h-2kx
 &=(v-\delta k^2b^2)[(w-\delta k)+\delta(u-k)]
       -\delta kb^2(u-k)(w-\delta k),\notag\\
 kC_mr_*&=q^2-yz.                                      \tag{9}
\end{align}
For any polynomial $S$ successive monic divisions give
\[
 S=S_c+(u-k)T_u+(w-\delta k)T_w
              +(v-\delta k^2b^2)T_v.
\]
If $S_c=0$, (9) yields the explicit polynomial lift
\begin{equation}
 C_mS=x(-uT_u-wT_w-b^2kwT_v)
      +y(\delta T_w+\delta kb^2T_v)+z(T_u+b^2wT_v).      \tag{10}
\end{equation}
Thus
\begin{equation}
 (A_m,B_m,C_m)\cap\mathfrak J\subset(x,y,z),\qquad
 [(A_m,B_m,C_m)^2+(D_0^2)]\cap\mathfrak J^2
                      \subset(x,y,z,h,q)^2.            \tag{11}
\end{equation}
The restriction to the unmixed generator $D_0^2$ is substantive.
For clarity, here is the complete constructive argument. In the first
inclusion write $F=R_AA_m+R_BB_m+R_CC_m$ and substitute
$A_m=y+kC_m$, $B_m=x+C_m$. The remaining coefficient of $C_m$ is
$S=kR_A+R_B+R_C$. The contact equality implies $S_c=0$ by cancellation
in the polynomial domain $\Rc$, so apply (10).

For the second inclusion write
$F=\sum_{i\le j}f_{ij}m_im_j+\rho D_0^2$, with
$(m_1,m_2,m_3)=(A_m,B_m,C_m)$. Substitute the same formulas and
$D_0=kC_m-h$ to obtain
\[
 F=Q_2+C_m(xH_x+yH_y+zH_z+hH_h)+C_m^2S,
 \quad Q_2\in(x,y,h)^2,\quad H_h=-2k\rho,\quad H_z=0.
\]
The contact value implies $S_c=0$. Absorb $C_m^2S$ into the four
$H$ coefficients by (10). Since $F\in\mathfrak J^2$, its first
derivatives vanish at contact, while
\[
 \det\left(\frac{\partial(x,y,z)}{\partial(u,w,v)}\right)_c
        =-\delta^4k^4b^4\ne0.
\]
The relation $r_*\in\mathfrak J^2$ has gradient coefficients
$(-2k,1,1,-1)$. Put
\[
 \lambda=2k\rho_c,\qquad J_i=H_i-\lambda(-2k,1,1,-1)_i.
\]
The last $J_i$ has zero contact
value; the determinant and the first derivative equations imply the
other three do too. This uses the adjugate and cancellation only to
prove zero polynomial identities; no determinant inverse is a lifting
coefficient. Finally
\[
 F=Q_2+2\rho_c(q^2-yz)+\sum_{i=1}^4r_i(C_mJ_i),
 \qquad(r_1,r_2,r_3,r_4)=(x,y,z,h),
\]
and each $C_mJ_i$ is lifted by (10). This proves (11), including all
coefficient-zero strata, in the unlocalized ring.

\subsection{The new positive-weight estimate}

Assign weights
\begin{equation}
 (\delta,k,b,a,t,u,w,v)=(0,2,6,0,0,2,0,16).             \tag{12}
\end{equation}
The zero-weight variables are bounded; every positive-weight variable
is infinitesimal in its own field. Positive weight of every polynomial
coefficient therefore makes it a bounded separated infinitesimal,
without any comparison of independent rates. A polynomial
$\Ow(s)$ has all monomials of weight at least $s$.

The short coquotients and (5) give
\begin{align*}
 C_i&=2K^6+\Ow(12)\quad(i=2,b,3,5,\alpha),&
 C_4&=2K^6+\Ow(16),& C_d&=4K^7+\Ow(12),\\
 Q&=K+\Ow(16),& T_n&=K^2+\Ow(20),& V_n&=K^2+\Ow(18),\\
 D_v&=1+\Ow(32),&Y_n&=\Ow(2),&B_0&=-K^3+\Ow(18),\\
 A_c&=2K^8+\Ow(12),&&\mathcal C=K^{12}+\Ow(16).
\end{align*}
The stronger $C_4$ estimate follows because every $\eta$-term in
$P_4$ also contains $\kappa$, and the clearing correction $Q-K$
has weight $16$. The two exact difference identities are
\[
 K(C_b-C_2)=\delta^2k^2Qb^2C_3,\qquad
 Q(C_3-C_5)=-k^2Kb^2C_b.
\]
Both differences therefore have weight at least $16$; since $K,Q$
have constant term one, this conclusion is polynomial and introduces
no inverse into a lifting coefficient.

The monomial weights of $(A_m,B_m,C_m,D_0)$ are $(16,16,14,18)$.
In the five lines of (6), keep the respective leading terms
\[
 8K^{14}D^2A_m^2,\quad 8K^{14}D^2B_m^2,\quad
 8K^{15}DC_m^2,\quad 8K^{15}DD_0^2,\quad
 -16K^{16}DB_mC_m.
\]
The first two errors have weight at least $44$. The third error has
weight at least $42$: its only possibly smaller cross term contains
$b^2B_mC_m$, of weight $42$. The fourth error has weight at least
$48$. In the fifth line the coefficient error on $B_mC_m$ has weight
at least $42$, and the separate $b^2B_m^2$ term has weight at least
$44$. Their sum is $8K^{14}DP+\Ow(42)$. Since $\wt(P)\geq28$
and $\mathcal C-K^{12}=\Ow(16)$, this proves
\[
 E_0=\Ow(42).
\]
Also $S_B,S_C=\Ow(12)$, and (7) gives
\begin{equation}
 L_1=\Ow(32),\qquad L_2=\Ow(26).                       \tag{13}
\end{equation}

It remains to check that the explicit conductor divisions preserve
strict positivity of coefficient weights. The contact substitutions
$u\mapsto k$, $w\mapsto\delta k$, $v\mapsto\delta k^2b^2$
never decrease weight; their weights change respectively as
$2\mapsto2$, $0\mapsto2$, $16\mapsto16$.
For $S=\Ow(s)$, its monic quotients obey
\[
 \wt(T_u)\geq s-2,\quad \wt(T_w)\geq s,\quad
 \wt(T_v)\geq s-16.
\]
The coefficients in (10) on $(x,y,z)$ consequently have weights
at least $(s-2,s-2,s-4)$.

For a first-normal polynomial of weight $p$, monomial division by
$(A_m,B_m,C_m)$ gives quotient weights at least
$(p-16,p-16,p-14)$. Hence $S=kR_A+R_B+R_C$ has weight at least
$p-16$, and the completed lift has coefficient weights at least
$p-20$. From (13) these bounds are $12$ for $L_1$ and $6$ for $L_2$.

For $E_0$ of weight at least $42$, its monomial quotients on
\[
 A_m^2,A_mB_m,A_mC_m,B_m^2,B_mC_m,C_m^2,D_0^2
\]
have weights at least $(10,10,12,10,12,14,6)$. Thus the initial
$Q_2$ coefficients have weight at least $6$, the initial
$(H_x,H_y,H_h)$ have bounds $(10,10,8)$, and $S$ has bound $10$.
The first use of (10) leaves bounds $(8,8,6,8)$ on
$(H_x,H_y,H_z,H_h)$. Since $\rho_c$ has weight at least $6$ and
$\lambda=2k\rho_c$ at least $8$, subtracting the gradient relation
leaves the same bounds on $(J_x,J_y,J_z,J_h)$. A second use of
(10) therefore leaves every scalar-lift coefficient with weight at
least $2$. The separate coefficient $2\rho_c$ of $q^2$ has bound
$6$. Every correction coefficient is thus strictly infinitesimal.

\subsection{The full strict Gram and restoration}

The eleven original soft roots are $g_0=x$, $g_4=uwv-\delta^2k^4b^2$,
the six differences $a_i^2-a_j^2$, and the three differences
$a_ia_j-a_\ell a_m$ over the three pairings of four indices.
They satisfy the exact identity
\begin{equation}
 P=g_0^2+g_4^2+\frac13\sum_{i<j}(a_i^2-a_j^2)^2
       +\frac23\sum_{\{ij,\ell m\}}(a_ia_j-a_\ell a_m)^2. \tag{14}
\end{equation}
The conductor roots $x,y,z,h,q$ are respectively
$g_0,a_1^2-a_3^2,a_2^2-a_3^2,a_3^2-a_4^2,a_1a_2-a_3a_4$.
Thus (8)--(13) express the complete scalar correction and both full
first normals as infinitesimal perturbations on this same finite
root frame. The leading diagonal in (14), multiplied by
$8K^2D\mathcal C$, has strictly positive residue.

There are no state-normal terms above degree two. To display the
entire hard block before multiplication by $\mathcal C$, put
\[
 C_B=H_3D_vh_s+H_5(1-h_s),
\]
\[
 B_Y=4RT\Delta^2\eta^2c_4DD_v
      +4RT\Delta^4k^6\eta^2c_\alpha D\,v^2
      +2RT\Delta^2\eta c_dD\,v^2.
\]
The entries on $(N,\delta M_B)$ are
\begin{align*}
 A_{11}&=H_2+H_bD_v+\Delta^2C_Bu^2v^2
                         +B_Yh_s^2u^2v^2/D^2,\\
 A_{12}&=-\frac{\Delta^2}{\delta}(H_3D_v-H_5)uv
                         -\frac{B_Y}{\delta D^2}h_suv,\\
 A_{22}&=\frac{\Delta^2}{\delta^2}(H_3D_v+H_5)
                         +\frac{B_Y}{\delta^2D^2}.
\end{align*}
Every apparent negative power of $\delta$ cancels: $\Delta/\delta=Q/K$
is a source unit, and $B_Y/\delta^2$ is bounded infinitesimal even
when $\delta$ is a unit. Hence this exact hard block has residue
$16I_2$. Its denominators are only positive pure-source or
pure-testing units. Together with (14) and the positive-weight
corrections, the full thirteen-root Gram for $\mathcal CG$ is bounded
and has strictly positive residue. Apply the admitted separated
strict-Gram lemma, restore all four packets $\mathcal C\Pi_i$,
and undo the positive source multiplier $\mathcal C\Delta^2\Lambda$.
This proves the entire original parent in the $Z$-unit chart with
bounded $W$, retaining both strict gaps and every original event.

\section{Every remaining testing chart}

\paragraph{Outer angle charts.}
The two preserved outer-angle identities require only
$d_{\rm B}/C,d_{\rm B}/D_{\rm e}$ infinitesimal and the harmonic
mass $\gamma=A[CD_{\rm e}/(C+D_{\rm e})]/d_{\rm B}^2$ bounded
with residue at least $1/4$. These hypotheses are independent of the
size of $e$. Indeed, put $B_c=2c_5-J\eta c_d$. The original source gives
\[
 D_{\rm e}=\frac{aB_c}{Jc_3},\quad
 \frac{d_{\rm B}}{D_{\rm e}}=\frac{2J\eta c_d}{B_c},\quad
 \frac{d_{\rm B}}C=\frac{a\eta c_d}{c_3},
\]
and
\[
 \gamma=\frac{c_4c_3B_c}{c_d^2(2Jc_3+aB_c)}.
\]
Here $\overline B_c=4$, $J+a=T$ has residue one, the denominator
has residue $64$, and the numerator has residue $16$. Thus both
strength ratios are infinitesimal and $\overline\gamma=1/4$.
The exact retained-current identities in
\texttt{c1\_uniform\_all\_tests.tex}, Lemmas 1--2, cover respectively
$1/v$ bounded and $v$ infinitesimal with $1/u$ bounded, with all
homogeneous infinity endpoints. Consequently only $u,v$
infinitesimal remain. The literal axes have a nonnegative core.

\paragraph{State charts.}
Normalize a nonzero state by a coordinate of maximal absolute value.
All its coordinates are then bounded and at least one is a unit. If
$X$ or $Y$ is a unit, the signed core has positive residue
$2(X^2+Y^2)$ and gives a direct separated square; restore all event
squares. Otherwise $X,Y$ are infinitesimal. If $Z$ is a unit, the
preceding section applies, including bounded unit $W$.

The only remaining case has $Z$ infinitesimal and $W$ a unit.
Normalize $W=1$ and put $Q_A=1+uY$, a pure-testing unit. Use
\[
 z=(r_2,B,\Delta r,\Delta r_5,\Delta\eta Y).
\]
In (2), retain the four nonnegative excess packets obtained from
$R/M=1+t/M$, $R/t=1+M/t$, $T/a=1+J/a$, $T/J=1+a/J$,
the complete $\alpha$ packet, and the positive $v^2Y^2$ part of the
last line. The remaining diagonal is
\[
 H_2z_1^2+H_bD_vz_2^2+H_3D_vz_3^2+H_5z_4^2
                   +4RTc_4DD_vz_5^2,
\]
and its complete negative correction is exactly
\[
 -4RTc_dD\frac{\Delta Z}{Q_A}z_5(z_1+z_2)
 -2RTc_dD\frac{\Delta^2\eta Z^2}{Q_A^2}(z_1+z_2)^2.
\]
This is the preserved $W$-unit identity, now on the larger domain:
$\Delta$ is bounded and $Z$ is infinitesimal, so both corrections
are infinitesimal. The diagonal residue is $8I_5$. The separated
strict-Gram lemma and restoration of all six packets prove this
last chart. There is no need for $\Delta$ itself to be infinitesimal.

The zero state is immediate. A complex state splits into its real and
imaginary quadratic evaluations because the parent matrix is real
symmetric. Every normalization and projective clearing factor is
pure testing. This completes the proof of Theorem~\ref{thm:bounded-delta}.

\section*{Provenance and focused exact check}

The source coquotients, full parent, retained packets, contact balances,
and normal formulas are those preserved in
\texttt{prior/prior\_core/c1/source\_definitions.tex}; the polynomial
conductor is from its adjacent \texttt{contact\_conductor\_proof.tex}.
The eleven-root identity and full hard block are in the preserved
\texttt{c1\_unbounded\_e\_z\_obligation.tex}. Their exact identities
do not require the old infinitesimal-$\delta$ hypothesis. The new
weight argument (12)--(13), the bounded-$W$ use, and the final
$W$-unit/$Z$-infinitesimal join are the scope extension proved here.
Exact paths and inherited hashes are in the companion routing ledger.

The adjacent standard-library checker reconstructs the source from
the seven short coquotients, performs only monic contact division and
exact removal of displayed monomial factors, and multiplies back the
complete scalar and both first-normal lifts. Its 33 assertions pass.
The actual minimum weights are $44,32,26$ on $(E_0,L_1,L_2)$ and
$12,14,12$ on the constructed correction families, stronger than the
human bounds above. No Gr\"obner basis, numerical positivity test,
parameter inversion, or prerecorded source expansion is used.
The checker corroborates the proof; it does not replace the displayed
conductor or the complete-parent chart join.
